\documentclass[a4paper,12pt]{jsarticle}
\usepackage[margin=25mm]{geometry}
\usepackage{amsmath,amssymb}
\usepackage{enumitem}

\begin{document}

\begin{center}
{\LARGE 数I・A 共通テスト本番レベル（図形）}
\end{center}

\vspace{1em}

\section*{問題}

以下では、図や会話、操作の説明を読み取りながら考える。必要に応じて、途中結果を文字式で整理してから計算すること。

\section*{第1問}
高校の探究活動で、2地点 $A,B$ から塔 $T$ の位置を測ることにした。  
平地上の2地点 $A,B$ の距離は $AB=80\text{ m}$ である。地点 $A$ から見た塔 $T$ への方位と、地点 $B$ から見た塔 $T$ への方位を測定したところ、平面図において
\[
\angle TAB=42^\circ,\quad \angle TBA=57^\circ
\]
であったとする（$T$ は線分 $AB$ と同じ側にある）。

さらに、塔の頂点を $P$、塔の足元を $T$ とし、地点 $A$ からの仰角を測ると
\[
\angle PAT=28^\circ
\]
であった。塔は地面に対して垂直であるとする。
必要なら
\[
\sin57^\circ=0.839,\ \sin81^\circ=0.988,\ \tan28^\circ=0.532,\ \tan27^\circ=0.510
\]
を用いてよい。

\begin{enumerate}[label=\arabic*.]
  \item 三角形 $ABT$ について、$\angle ATB$ を求めよ。
  \item $AT$ を求めよ。
  \item 2.で求めた結果を用いて、塔の高さ $PT$ を求めよ。
  \item 測定誤差により $\angle PAT$ が $28^\circ$ ではなく $27^\circ$ だったと仮定すると、塔の高さは何 m 変化するか（増減も含めて）求めよ。
\end{enumerate}

\section*{第2問}
半径 $6$ の円 $O$ があり、円周上に2点 $A,B$ をとる。  
弦 $AB$ の長さは $6\sqrt3$ である。点 $A$ における接線上に点 $P$ をとり、$\triangle PAB$ を考える。  
また、点 $P$ は「$PA=6$」を満たすようにとる。

\begin{enumerate}[label=\arabic*.]
  \item 中心角 $\angle AOB$（小さい方）を求めよ。
  \item 弧 $AB$（小さい方）の長さと、扇形 $AOB$（小さい方）の面積を求めよ。
  \item $\angle PAB$ を求めよ。
  \item 三角形 $PAB$ の面積を求めよ。
  \item 線分 $PB$ の長さを求めよ。
\end{enumerate}

\section*{第3問}
平行四辺形 $ABCD$ において
\[
AB=8,\quad AD=5,\quad \angle DAB=60^\circ
\]
とする。対角線 $AC,BD$ の交点を $O$ とする。  
点 $H$ を、点 $A$ から直線 $BD$ に下ろした垂線の足とする。

\begin{enumerate}[label=\arabic*.]
  \item 平行四辺形 $ABCD$ の面積を求めよ。
  \item 対角線 $BD$ の長さを求めよ。
  \item 三角形 $AOD$ の面積を求めよ。
  \item 線分 $AH$ の長さを求めよ。
  \item $BH:HD$ を求めよ。
\end{enumerate}

\section*{第4問}
直方体 $ABCD\text{-}EFGH$ において
\[
AB=6,\quad BC=4,\quad AE=8
\]
とする（$ABCD$ が下底面、$EFGH$ が上底面）。  
点 $M$ は辺 $CG$ の中点、点 $N$ は辺 $AE$ を $AN:NE=1:3$ に内分する点とする。  
平面 $\alpha$ を、3点 $B,M,N$ を通る平面とする。

\begin{enumerate}[label=\arabic*.]
  \item 座標を
  \[
  A(0,0,0),\ B(6,0,0),\ C(6,4,0),\ E(0,0,8)
  \]
  とおいたとき、点 $M,N$ の座標を求めよ。
  \item 平面 $\alpha$ の方程式を求めよ。
  \item 点 $D$ から平面 $\alpha$ までの距離を求めよ。
  \item 線分 $AG$ と平面 $\alpha$ の交点を $X$ とするとき、$AX:XG$ を求めよ。
\end{enumerate}

\end{document}
